
\section{Vorbereitung}
\subsection{Aufgabe 1 - Temperaturabhängige Widerstände}
Um eine Temperatur in eine Elektrische Größe umzuwandeln, können temperaturabhängige Widerstände, auch Thermistor genannt, verwendet werden.
Das sind Widerstände, deren Widerstandswert sich mit der Temperatur ändert. In Schaltung \ref{fig:vorbereitung_1} ist eine Messschaltung dargestellt, mit dem man eine Messung des Widerstandes $R_T$ durchführen kann.
Hierbei wird dazu ein Spannungsteiler verwendet, der aus einem festen Widerstand $R_0$ und dem temperaturabhängigen Widerstand $R_T$ besteht.

\begin{Schaltung}[H]
    \centering
    \begin{circuitikz}[european]
        % Coordinates: Starting from (0, 3) for the top of R0

        % Top Input Terminal (extended to the left)
        \draw (-2,3) to[short, o-] (0,3);

        % Resistor R0
        \draw (0,3) to[R, -, name=R0, label=$R_0$] (0,1.5) coordinate (mid);
        %\node[right=2pt] at (0,2.25) {$R_0$}; % R0 label placement

        % Resistor RT (Variable Resistor)
        \draw (mid) to[thermistor, -, name=RT, label=$R_T$, invert] (0,0);

        % Bottom Input Terminal (extended to the left)
        \draw (-2,0) to[short, o-] (0,0);

        % Output Terminals
        % Mid Connection Point
        \node[black, fill=black, inner sep=1.5pt, circle] at (mid) {};
        \draw (mid) to[short, -o] (2,1.5); % Top output terminal

        % Bottom Connection Point
        \node[black, fill=black, inner sep=1.5pt, circle] at (0,0) {};
        \draw (0,0) to[short, -o] (2,0);   % Bottom output terminal

        % U_0 Voltage Arrow (Input - specifically placed to the left)
        \draw[->] (-1.5,2.5) -- (-1.5,0.5) node[pos=0.5, left] {$U_0$};

        % U_T Voltage Arrow (Output - specifically placed to the right, pointing down)
        \draw[->] (2.5,1.25) -- (2.5,0.25) node[pos=0.5, right] {$U_T$};

    \end{circuitikz}
    \caption{Messschaltung für temperaturabhängigen Widerstand $R_T$}
    \label{fig:vorbereitung_1}
\end{Schaltung}

Hierzu muss $R_0$ korrekt dimensioniert werden.
Es gilt den folgenden Spannungsteiler:
\begin{equation}
    U_T = U_0 \cdot \frac{R_T}{R_0 + R_T}
\end{equation}
Jetzt leiten wir $U_T$ nach $R_T$ ab, um die temperaturabhängige Spannungsänderung zu bestimmen:
\begin{equation}
    \label{eqn:temp_spannungsaenderung}
    \frac{\partial U_T}{\partial R_T} = \frac{(R_T+R_0)U_0-U_0R_T}{(R_T+R_0)^2} = U_0 \frac{R_0}{(R_T+R_0)^2}
\end{equation}
Um den Wert von $R_0$ zu ermitteln, bei dem die temperaturabhängige Spannungsänderung maximal ist, leiten wir die Gleichung \ref{eqn:temp_spannungsaenderung} nach $R_0$ ab und setzen diese gleich Null:
\begin{equation}
    \frac{\partial}{\partial R_0}\left( U_0 \frac{R_0}{(R_T+R_0)^2} \right) = U_0 \frac{R_T-R_0}{(R_T+R_0)^3}
\end{equation}
\begin{align}
    \Rightarrow U_0 \frac{R_T-R_0}{(R_T+R_0)^3} & = 0           \\
    R_T - R_0                                   & = 0 \nonumber \\
    \boxed{R_0 = R_T} \nonumber
\end{align}

Daraus folgt, dass die maximale temperaturabhängige Spannungsänderung sich bei einem Widerstandswert von $R_0 = R_T$ ergibt.

\subsection{Aufgabe 2 - PTC, NTC}
\subsubsection{pt100 Widerstand}
Der Pt100 ist ein Platin-Widerstand, welches bei seiner Bezugstemperatur (im Labor \SI{0}{\degreeCelsius}) einen Widerstand von \SI{100}{\ohm} aufweist.
Um nun den Widerstand bei einer Temperaturänderung von 1K zu bestimmen, verwenden wir die Formel eines PTC-Widerstands.
\begin{equation}
    R_{\vartheta} = R_0 \left[
        1
        + \num{3.90802e-3} \, \frac{\vartheta - \vartheta_0}{\mathrm{K}}
        - \num{5.80195e-7} \, \frac{(\vartheta - \vartheta_0)^2}{\mathrm{K^2}}
        \right]
\end{equation}

Bei PTC-Widerständen kann man in dem Temperaturbereich von \SI{0}{\degreeCelsius} bis \SI{850}{\degreeCelsius} die Kennlinie linearisieren.
\begin{figure}[H]
    \includegraphics[width=\linewidth]{vorbereitung_kennlinie_pt100.png}
\end{figure}
Unter dieser Linearisierung ergibt sich eine Widerstandsänderung pro Kelvin von $\SI{0.375}{\frac{\ohm}{\kelvin}}$. \\
Daraus folgt: \\
$R_{25} = \SI{100}{\ohm}$ \\
$R_{26} = \SI{100.375}{\ohm}$ \\

Die relative Spannungsänderung $\frac{\Delta U_T}{U_0}$ ergibt sich damit zu:
\begin{align}
    \frac{\Delta U_T}{U_0} & = \frac{R_{T2}}{R_0 + R_{T2}} - \frac{R_{T1}}{R_0 + R_{T1}}          \nonumber                                            \\
                           & = \frac{\SI{100.375}{\ohm}}{\SI{100}{\ohm} + \SI{100.375}{\ohm}} - \frac{\SI{100}{\ohm}}{\SI{100}{\ohm} + \SI{100}{\ohm}} \\
                           & = \num{935.745e-6} \nonumber                                                                                              \\
                           & = \SI{0.09}{\percent}
\end{align}

\subsubsection{Thermistor Typ B57863S863}
Unser Thermistor ist ein NTC-Widerstand, dessen Widerstand mit steigender Temperatur abnimmt.
In diesem Fall ist die Materialkonstante $B = \SI{3988}{\kelvin}$ und der Nennwiderstand $R_N = \SI{10}{\kilo\ohm}$ bei einer Nenntemperatur von $\vartheta = \SI{25}{\degreeCelsius}$. \\
Allgemein gilt für NTC-Widerstände die folgende Formel:
\begin{equation}
    R_T = R_N \cdot e^{B \left( \frac{1}{T} - \frac{1}{T_N} \right)}
\end{equation}
Damit lassen sich auch die Widerstände für die Temperaturen \SI{25}{\degreeCelsius} und \SI{26}{\degreeCelsius} berechnen:
\begin{align*}
    R_{26}             & = \SI{10}{\kilo\ohm} \cdot e^{3988 \left( \frac{1}{299.15 \, \mathrm{K}} - \frac{1}{298.15 \, \mathrm{K}} \right)} = \SI{9.562}{\kilo\ohm}
    \\
    \frac{U_{26}}{U_0} & = \frac{R_{26}}{R_{26} + R_0} = \frac{\SI{9.562}{\kilo\ohm}}{\SI{9.562}{\kilo\ohm} + \SI{10}{\kilo\ohm}} = \num{0.4888} = \SI{48.88}{\percent} \\
    \\
    \frac{U_{25}}{U_0} & = \frac{R_{25}}{R_{25} + R_0} = \frac{\SI{10}{\kilo\ohm}}{\SI{10}{\kilo\ohm} + \SI{10}{\kilo\ohm}} = \num{0.5} = \SI{50}{\percent}             \\
\end{align*}
Die relative Spannungsänderung $\frac{\Delta U_T}{U_0}$ ergibt sich damit zu:
\begin{align*}
    \frac{\Delta U_T}{U_0} & = \frac{U_{26} - U_{25}}{U_0} = \num{0.4888} - \num{0.5} = -\num{11.19e-3} = -\SI{1.119}{\percent}
\end{align*}

\subsection{Aufgabe 3 - Auf- und Entladevorgang einer RC-Kombination}
Aufladevorgang:
\begin{equation}
    U_C(t) = U_0 \left( 1 - e^{-\frac{t}{RC}} \right)
\end{equation}
\\
Entladevorgang:
\begin{equation}
    U_C(t) = U_0 \cdot e^{-\frac{t}{RC}}
\end{equation}

\subsection{Aufgabe 4 - Grafische Bestimmung von der Zeitkonstante}
\label{sec:vorbereitung_aufgabe4}
Zur Bestimmung der Zeitkonstante $\tau$ aus der Kurve des Spannungsverlaufs $U(t)$ stehen mehrere grafische Methoden zur Verfügung.
Die erste Möglichkeit ist die Bestimmung des Wertes von $t$ bei dem die Kurve von $U(t)$ 63\% des Endwertes erreicht hat. \\
Die zweite Möglichkeit ist die Bestimmung des Wertes von $t$ bei dem die Tangente der Kurve an der Stelle $t=t_0$ die Asymptote schneidet. \\
Beide Möglichkeiten liefern den gleichen Wert für die Zeitkonstante $\tau$.

\subsection{Aufgabe 5 - Verlustleitung bei Wärmeabfuhr bei Halbleiterbauelementen}
Zur mathematischen Berechnung der Wärmeabfuhr bei Halbleiterbauelementen können wir das resultierende thermische Netzwerk als Analog zu einem elektrischen Netzwerk betrachten mit folgenden Äquivalenzen:
\begin{align*}
    \text{Verlustleistung } P_V                  & \quad \hat{=} \quad \text{Elektrischer Strom } I      \\
    \text{Temperaturdifferenz } \Delta \vartheta & \quad \hat{=} \quad \text{Elektrische Spannung } U    \\
    \text{Temperatur } \vartheta                 & \quad \hat{=} \quad \text{Knotenpotenzial } U         \\
    \text{Wärmewiderstand } R_{th}               & \quad \hat{=} \quad \text{Elektrischer Widerstand } R \\
    \text{Wärmekapazität } C_{th}                & \quad \hat{=} \quad \text{Elektrische Kapazität } C
\end{align*}
Daraus kann man bekannte Gesetze der Elektrotechnik auf die Wärmeabfuhr übertragen.
\begin{equation}
    \label{eqn:Ohmsches_Gesetz_Waerme}
    \Delta \vartheta = P_V \cdot R_{th} \quad\text{(Ohmsches Gesetz)}
\end{equation}
Bei der Wärmeabfuhr erfolgt von der Sperrschicht über das Gehäuse und den Kühlkörper zur Umgebung.
Die Verlustleistung fließt somit durch eine Reihenschaltung von Wärmewiderständen.
\begin{Schaltung}[H]
    \begin{circuitikz}[european]

        % Koordinaten definieren für saubereren Code
        \coordinate (Start) at (0,0);
        \coordinate (NodeC) at (4,0);
        \coordinate (NodeS) at (8,0);
        \coordinate (End) at (11,0);

        % Haupfad mit Widerständen und Strompfeil
        \draw (Start)
        to[short, i=$P_v$] (1.5,0)
        to[R, l_=$R_{\text{th,JC}}$] (NodeC)
        to[R, l_=$R_{\text{th,CS}}$] (NodeS)
        to[R, l_=$R_{\text{th,S}}$] (End);

        % Start-Klemme (kleiner Kreis)
        \node[ocirc] at (Start) {};

        % Vertikale Linien für die Temperaturen
        % Wir zeichnen Linien nach unten von den Knotenpunkten
        \draw ($(Start)+(0,-0.2)$) -- ++(0,-1.5) node[below] {$\theta_J$};
        \draw ($(NodeC)+(0,-0.2)$) -- ++(0,-1.5) node[below] {$\theta_C$};
        \draw ($(NodeS)+(0,-0.2)$) -- ++(0,-1.5) node[below] {$\theta_S$};
        \draw ($(End)+(0,-0.2)$)   -- ++(0,-1.5) node[below] {$\theta_a$};

        % Die Wand / Umgebung (Schraffur am Ende)
        % Vertikale Linie
        \draw (End) ++(0,1) -- ++(0,-2.5);

        % Schraffur-Linien (Loop für den Look)
        \foreach \y in {1, 0.6, ..., -1.4} {
                \draw (End) ++(0,\y) -- ++(0.5, 0.5);
            }
    \end{circuitikz}
    \caption{Ersatzschaltbild Thermisches Netzwerk zur Wärmeabfuhr bei Halbleiterbauelementen}
    \label{fig:thermisches_netzwerk}
\end{Schaltung}
Um die Verlustleistung $P_V$ bei einer Temperatur $\vartheta = \SI{125}{\degreeCelsius}$ zu ermitteln muss man
also nur die Formel \ref{eqn:Ohmsches_Gesetz_Waerme} nach $P_V$ umstellen und die entsprechenden Wärmewiderstände addieren.
\begin{equation}
    P_V = \frac{\Delta \vartheta}{R_{\text{th,gesamt}}}
\end{equation}
Die Wärmewiderstände setzen sich dabei wie folgt zusammen:
\begin{align*}
    R_{\text{th,gesamt}}     & = R_{\text{th,JC}} + R_{\text{th,CS}} + R_{\text{th,SA}}
    = \SI{3.13}{\kelvin\per\watt} + \SI{1}{\kelvin\per\watt} + \SI{15}{\kelvin\per\watt}    \\
                             & = \SI{19.13}{\kelvin\per\watt}                               \\
    R_{\text{th,gesamt,iso}} & = R_{\text{th,JC}} + R_{\text{th,CS,iso}} + R_{\text{th,SA}}
    = \SI{3.13}{\kelvin\per\watt} + \SI{2}{\kelvin\per\watt} + \SI{15}{\kelvin\per\watt}    \\
                             & = \SI{20.13}{\kelvin\per\watt}
\end{align*}
Ohne eine Isolierscheibe zwischen Gehäuse und Kühlkörper ergibt sich für die Verlustleistung:
\begin{equation}
    P_V = \frac{\Delta \vartheta_J}{R_{\text{th,JC}} + R_{\text{th,CS}} + R_{\text{th,SA}} }
    = \frac{\SI{125}{\kelvin} - \SI{25}{\kelvin}}{\SI{19.13}{\kelvin\per\watt}}
    = \SI{5.23}{\watt}
\end{equation}
Mit einer Isolierscheibe zwischen Gehäuse und Kühlkörper ergibt sich für die Verlustleistung:
\begin{equation}
    P_{V, \text{iso}} = \frac{\Delta \vartheta_J}{R_{\text{th,JC}} + R_{\text{th,CS,iso}} + R_{\text{th,SA}} }
    = \frac{\SI{125}{\kelvin} - \SI{25}{\kelvin}}{\SI{20.13}{\kelvin\per\watt}}
    = \SI{4.97}{\watt}
\end{equation}

\subsection{Aufgabe 6 - Wärmekapazität bei Wärmeabfuhr bei Halbleiterbauelementen}
Die überschlägige Berechnung der spezifischen Wärmekapazität $c_K$ des im Versuch verwendeten Kühlkörpers erfolgt über folgende Formel:
\begin{equation}
    c_K = m \cdot c_{Al} \cdot \Delta T
    = \SI{3.3}{\gram} \cdot \SI{0.9}{\frac{\kilo\joule}{\kilogram\kelvin}}
    = \SI{2.97}{\joule\per\kelvin}
\end{equation}

\subsection{Aufgabe 7 - Stationäre Temperatur}
Um die Zeit zu ermitteln, die der Kühlkörper benötigt, um eine stationäre Temperatur zu erreichen, modellieren wir das thermische Netzwerk als RC-Glied.
Die Zeitkonstante $\tau$ ergibt sich aus dem Produkt von Wärmewiderstand und Wärmekapazität:
\begin{align}
    \tau              & = R_{\text{th,gesamt}} \cdot C_K = \SI{19.13}{\kelvin\per\watt} \cdot \SI{2.97}{\joule\per\kelvin}     \\ & = \SI{56.8}{\second} \nonumber \\
    \tau_{\text{iso}} & = R_{\text{th,gesamt,iso}} \cdot C_K = \SI{20.13}{\kelvin\per\watt} \cdot \SI{2.97}{\joule\per\kelvin} \\ & = \SI{59.8}{\second} \nonumber
\end{align}
Die Temperatur erreicht nach $5\tau$ ca. 99\% des Endwertes, wird also praktisch stationär.
Somit ergibt sich die Zeit bis zur stationären Temperatur zu:
\begin{align}
    t_{\text{stationär}}     & = 5 \cdot \tau = 5 \cdot \SI{56.8}{\second} = \SI{284.05}{\second} = \SI{4.73}{\minute}             \\
    t_{\text{stationär,iso}} & = 5 \cdot \tau_{\text{iso}} = 5 \cdot \SI{59.8}{\second} = \SI{299.9}{\second} = \SI{4.98}{\minute}
\end{align}

\subsection{Aufgabe 8 - Berechnung von RT}
Um den Widerstand $R_T$ in Abhängigkeit von $U_{ref}$ und $R_0$ zu berechnen, wird der Spannungsteiler aus Schaltung \ref{fig:vorbereitung_1} verwendet.
\begin{alignat}{2}
    \label{eqn:a8}
    % Zeile 1: Kein Pfeil, leere erste Spalten
                    & \quad & U_{T}                                            & = U_{ref} \cdot \frac{R_T}{R_0 + R_T} \nonumber   \\
    % Zeile 2: Pfeil in Spalte 1, Gleichung in Spalte 2
    \Leftrightarrow &       & R_T                                              & = (R_T + R_0) \cdot \frac{U_T}{U_{ref}} \nonumber \\
    \Leftrightarrow &       & R_T - R_T \cdot \frac{U_T}{U_{ref}}              & = R_0 \cdot \frac{U_T}{U_{ref}} \nonumber         \\
    \Leftrightarrow &       & R_T \cdot \left( 1 - \frac{U_T}{U_{ref}} \right) & = R_0 \cdot \frac{U_T}{U_{ref}}                   \\
    \Leftrightarrow &       & R_T \cdot \frac{U_{ref} - U_T}{U_{ref}}          & = R_0 \cdot \frac{U_T}{U_{ref}} \nonumber         \\
    \Leftrightarrow &       & R_T                                              & = R_0 \cdot \frac{U_T}{U_{ref} - U_T} \nonumber
\end{alignat}

\subsection{Aufgabe 9}
Die Spannung $U_T$ wird vom SAR-ADC gemessen und als Zahlenwert $Z$ ausgegeben. Allgemein gilt für $Z$, der bei einem 12-Bit SAR-ADC mit einer Eingangsspannung $U_T$ und einer Vollausschlagspannung $U_{\text{SAR,max}} = 2U_{\text{ref}}$ und $Z_{\text{SAR,max}} = 4095 = 2^{12} - 1$:
\begin{alignat}{2}
    \label{eqn:a9}
                    & \quad & \frac{Z}{Z_{\text{SAR,max}}} & = \frac{U_T}{U_{\text{SAR,max}}}                                 \\
    \Leftrightarrow &       & Z                            & = \frac{U_T}{2U_{\text{ref}}} \cdot Z_{\text{SAR,max}} \nonumber
\end{alignat}

\subsection{Aufgabe 10}
Unter Verwendung der Ergebnisse aus Aufgabe 8 und 9 kann nun gezeigt werden, dass
$R_T$ unabhängig von $U_{\text{ref}}$ ist.
Dazu setzen wir die Gleichung \ref{eqn:a9} in die Gleichung \ref{eqn:a8} ein:
\begin{alignat}{2}
    \label{eqn:a10}
                    & \quad & R_T & = R_0 \cdot \frac{\frac{Z}{Z_{\text{SAR,max}}} \cdot 2U_{\text{ref}}}{U_{\text{ref}} - \frac{Z}{Z_{\text{SAR,max}}} \cdot 2U_{\text{ref}}} \nonumber \\
    \Leftrightarrow &       &     & = R_0 \cdot \frac{\frac{2Z}{Z_{\text{SAR,max}}} \cdot U_{\text{ref}}}{U_{\text{ref}} \left( 1 - \frac{2Z}{Z_{\text{SAR,max}}} \right)}               \\
    \Leftrightarrow &       &     & = R_0 \cdot \frac{2Z}{Z_{\text{SAR,max}} - 2Z} \nonumber
\end{alignat}

\subsection{Aufgabe 11}
Unter der Annahme, dass der SAR-ADC den Zahlenwert $Z = 1024$ ausgibt und $R_0 = \SI{10}{\kilo\ohm}$, ergibt sich für den Widerstand $R_T$:
\begin{equation}
    R_T = R_0 \cdot \frac{2Z}{Z_{\text{SAR,max}} - 2Z} = \SI{10}{\kilo\ohm} \cdot \frac{2 \cdot 1024}{4095 - 2 \cdot 1024} = \SI{10.00489}{\kilo\ohm}
\end{equation}
Um die zugehörige Temperatur zu bestimmen wird die Steinhart-Hart-Gleichung verwendet:
\begin{equation}
    \frac{1}{T} = A + B \cdot \ln(R_T) + C \cdot (\ln(R_T))^3
\end{equation}
Mit den Konstanten aus der Vorbereitung
\begin{align*}
    A & = \SI{1.12559726e-03}{\per\kelvin} \\
    B & = \SI{2.34801612e-04}{\per\kelvin} \\
    C & = \SI{8.63986114e-08}{\per\kelvin}
\end{align*}
und $R_T = \SI{10.00489}{\kilo\ohm}$ ergibt sich für die Temperatur:
\begin{align*}
    \frac{1}{T} & = \SI{1.126e-3}{\per\kelvin} + \SI{2.348e-4}{\per\kelvin} \cdot \ln\left(\frac{\SI{10004.89}{\ohm}}{\SI{1}{\ohm}}\right) \\
                & \quad + \SI{8.640e-8}{\per\kelvin} \cdot \ln^3\left(\frac{\SI{10004.89}{\ohm}}{\SI{1}{\ohm}}\right)                      \\
                & = \SI{3.3558e-3}{\per\kelvin}                                                                                            \\
\end{align*}

\begin{equation}
    \Rightarrow \boxed{T = \SI{297.989}{\kelvin} \approx \SI{25}{\degreeCelsius}}
\end{equation}


